\lecture{1}{Thu 25 Aug 2022 12:01}{Properties of Complex Numbers}

\section{Review}
\subsection{Complex Numbers}
\subsubsection{Basic Properties}
\begin{itemize}
	\item $e^{x+iy } = e^x \left( \cos y + i \sin y \right)$ 
	\item $e^0 = 1$
	\item $e^{2\pi i} = 1$
	\item $e^{z + 2\pi i} = e^z$ for every $z$. Smallest period $2 \pi i$
	\item $e^z \neq  0$
	\item $e^{-z } = \frac{1}{e^z}$
	\item $e^{i \pi} = -1$
	\item $e^{i \frac{\pi}{2}} = i$
	\item $e^{-i\frac{\pi}{2}} = -i$
	
\end{itemize}

Geometry.
\begin{enumerate}
	\item $|e^{i \theta}| = 1$
\item  $\{z: |z| = 1\} $ is the unit circle.
\end{enumerate}

Argument.
\begin{itemize}
	\item (algebraic form) $z = x+iy$
	\item (exponential) $z = |z| e^{i \text{arg} z} = |z| e^{i \theta}$
\end{itemize}

Algebraic form is good for addition, exponential form for multiplication.

\subsubsection{ Complex exponentiation and applications } 
\begin{itemize}
	\item 
		\begin{theorem}[De Moivre]
			$\left( \cos \theta + i \sin \theta \right) ^{n} = e^{i n \theta} = \cos n \theta + i \sin n \theta$, $n$ integer
		\end{theorem}
	\item 	\begin{example}[Application to trignometric identities]
			\begin{align*}
			\cos 3 \theta &= \text{Re}\left( \cos \theta + i \sin \theta \right) ^{3}\\
				      &= \text{Re}\cos^3 \theta + 3 i \cos^2 \theta \sin \theta - 3 \cos \theta \sin ^2 \theta - i \sin ^3 \theta\\
				      &= \cos^3 \theta - 3 \cos \theta \sin^2 \theta
		\end{align*}
		\end{example}
		
	\item  Definition of trignometric functions in term of complex exponentiation
		\begin{equation*}
			\begin{cases}
				\cos \theta = \frac{1}{2}(e^{i \theta} + e ^{- i \theta}) \\
				\sin \theta = \frac{1}{2i}(e^{i \theta} - e ^{- i \theta}) 
			\end{cases}
		\end{equation*}
We may define complex values of trig functions by extending $\theta$ from real number to complex number.
\end{itemize}


\begin{example}
	Solve $z^n = 1$.
\begin{proof}[Solution]
\begin{align*}
	z &= r e^{i \theta}\\
	z^n &= r^n e^{i n \theta} = e^{2 \pi i k}\\
	\intertext{$r = 1$ by comparing abs values.}
	i n \theta &= 2 \pi i k\\
	\theta &= \frac{2 \pi k }{n }, k \in  \mathbb{Z}\\
	z_k &= e^{i \theta_k} = e^{\frac{2 \pi i }{n}k}, k = 0,1,\ldots,n-1
.\end{align*}
We have $n $ distinct solutions. 
\end{proof}
\end{example} 
\begin{example}
Solve the equation $z^n = w$ where $w$ is a given complex number.
\begin{proof}[Solution]
\begin{align*}
	r^n e ^{i n \theta} &= \rho e^{i \phi}\\
	r &= \sqrt[n]{\rho} \\
	i n \theta &= i \varphi + 2 \pi i k\\
	\theta_k &= \frac{\varphi}{n }+\frac{2 \pi k }{n}\\
	z_k &= \rho^{\frac{1}{n } } e^{i \theta_k}
.\end{align*}
\end{proof}
\end{example}
